\section{Interventions and Audit Contracts}
\label{sec:intervention}
All arms replay the same frozen endpoint with the same evidence package, model configuration, output schema, and family scorer. We intervene only on the procedure exposure or on where an already computed procedure output enters the harness. The scientific unit is the endpoint; repeated executions are robustness observations rather than additional tasks.

\subsection{Base N/G/T stress test}
N adds no experimental helper. G exposes a frozen target-blind \emph{stress helper} using the same callable surface and matched implementation complexity as T, but it is forbidden from performing the registered temporal operation. G is not presented as a benign, optimal, or representative generic assistant: its purpose is to stress the common T-versus-generic evaluation pattern and reveal when control-arm degradation changes the verdict. For cutoff validity, G can validate package structure but cannot compare release dates with the cutoff; for release alignment, it cannot map publication dates to reporting periods; for grounding, it cannot perform the mechanism-specific contextual selector. Targeted and G source implementations are matched in lexical-token and AST-node count for each family.

This base test is useful precisely because G need not be benign. T--G is a stress contrast against structured but mechanism-misaligned assistance, while T--N determines whether anything was actually repaired. We do not use T--G alone to estimate a generic-helper treatment effect. The DeepSeek cutoff false positive in Section~\ref{sec:results} shows why the N anchor is mandatory whenever such a stress comparator is used.

\subsection{Same-surface no-operation placebo G$_0$}
G$_0$ executes the same assigned-helper path as T but its frozen function is simply \texttt{def skill(package, context): return \{\}}. It reads neither evidence nor task context. Consequently T and G$_0$ differ in whether the target operation output is present while sharing the same helper-exposure surface. We do not label G$_0$ behaviorally neutral at the stratum level. Its empirical difference from N is reported directly.

The final audit therefore does not combine T--N and T--G$_0$ into a single pure-operation estimand. T--N carries the \emph{net-repair} claim. T--G$_0$ is a \emph{same-surface operation-output contrast}; G$_0$--N reports how much the placebo/surface itself perturbs N. Their observed means satisfy Eq.~\ref{eq:decomp}. When G$_0$ harms or helps N, the T--G$_0$ contrast is conditional on that perturbed surface and is reported as such. The separately frozen A1/A2 G$_0$ studies remain diagnostics of wrapper perturbation; their full-track equivalence result is not transported into every small load-bearing stratum.

\subsection{Context-materialization surface control R$_{\mathrm{surf}}$}
R$_{\mathrm{surf}}$ isolates a narrower implementation question. It runs the exact targeted-operation source on the same candidates and computes the exact same operation output as T. Instead of exposing that output as \texttt{reusable\_helper\_output}, it adds the object as \texttt{retrieval\_operation\_output} inside the ordinary evidence context and exposes no callable helper. A zero-call audit verifies for every tested endpoint that removing this one added field reconstructs the original evidence package, operation-output hashes match T exactly, and candidate evidence is neither removed nor reordered.

Because R$_{\mathrm{surf}}$ deliberately reuses T's operation output, it is a \emph{surface-isolation control}, not an independent temporal-retrieval algorithm. It asks whether placing the same precomputed output on the T surface rather than in ordinary context changes the final answer under this forced one-answer harness. It does not test adaptive callable use during generation, whether a learned temporal retriever can discover or approximate the operation, or whether persistent callable state has an advantage in multi-turn deployment.

\subsection{Frozen outcomes, thresholds, and evidence layers}
Before each source-native execution tranche, endpoint identity, family outcome, repeat count, condition order, model configuration, and procedure source hashes are frozen. A cutoff success requires a nonempty cited set and latest selected reference that are all available by the target cutoff. A release-alignment success requires the correct reporting-period identity and aligned value. A grounding success requires the requested evidence slots to be filled by frozen first-party sentence IDs in the correct order. Auxiliary stricter scorers are reported only as sensitivity analyses.

The original R$_{\mathrm{surf}}$ contract freezes a 10-point material effect threshold on the 18-endpoint DeepSeek portfolio, before R$_{\mathrm{surf}}$ outcomes. The margin is conservative relative to the already observed 20--60 point scale of the load-bearing targeted effects available when the surface experiment was designed. We report the preregistered 90\% endpoint-bootstrap interval, a post-hoc 95\% bootstrap sensitivity interval, and a paired-endpoint TOST at the unchanged $\pm10$-point margin. The TOST is a statistical sensitivity analysis of the same frozen decision boundary rather than a replacement gate. For G$_0$, A1 and A2 remain unpooled; exact-stratum G$_0$--N diagnostics show where full-track equivalence does and does not transport.

The initial BEA/NOAA replay freezes pre-April and April--May periods before outcomes. A separate post-hoc EIA compatibility probe uses 16 fixed weekly releases to construct 12 deterministic cutoff endpoints: source indices 3--14, with four releases before April 1 assigned to C3-R4 and eight later releases assigned to C4-R4. Endpoint construction, scorer, T1/G1 hashes, and randomized condition order were frozen before EIA model outcomes. However, EIA was selected post-review because its dated weekly releases are structurally compatible with the cutoff mechanism, so C4-R4 is an illustrative mechanism-compatibility probe rather than prospective cross-domain confirmation.
